\chapter{Rayleigh分布}
    \section{定義}
        \begin{shadebox}
            $\sigma>0$とする。
            母数$\sigma$のRayleigh分布の連続確率密度関数を次式で定義する。
            \[ f(x;\sigma) := \frac{x}{\sigma^2}\exp\left(-\frac{x^2}{2\sigma^2}\right) \]
        \end{shadebox}
    \section{正規分布する成分を持つ2次元ベクトルのノルムとの関係}
        \begin{shadebox}
            $\bm{X}\in\realNumbers^2$の各成分が正規分布$N(0,\sigma^2)$に従う独立な確率変数であるとき、$\norm{\bm{X}}_2$は母数$\sigma$のRayleigh分布に従う。
        \end{shadebox}
        \begin{proof}
            \quad\par
            $N(0,\sigma^2)$の確率密度関数を$p$と表す。
            $r\geq 0$とする。
            $\Omega(r)$を半径$r$の円盤領域とする。
            \begin{align*}
                \Pr{\norm{\bm{X}}_2 \leq r} &= \int_{\Omega(r)} p(x_1)p(x_2) \mathrm{d}x_1\mathrm{d}x_2 = \integrate{0}{r}{\integrate{0}{2\pi}{p(r'\cos\theta)p(r'\sin\theta)r'}{}{\theta}}{}{r'} \\
                &= \integrate{0}{r}{\integrate{0}{2\pi}{\frac{r'}{2\pi\sigma^2}\exp\left(-\frac{r'^2}{2\sigma^2}\right)r'}{}{\theta}}{}{r'} = \integrate{0}{r}{\frac{r'}{\sigma^2}\exp\left(-\frac{r'^2}{2\sigma^2}\right)}{}{r'}
            \end{align*}
            よって$\norm{\bm{X}}_2$の確率密度関数は次式である。
            \[ \derivLong{\Pr{\norm{\bm{X}}_2 \leq r}}{r}{} = \frac{r}{\sigma^2}\exp\left(-\frac{r^2}{2\sigma^2}\right) \]
        \end{proof}
